% Auto-generated by scripts/06_emit_structure.py
% Source: scans 178–181
% Section id: appendix/2
% Phase 3 deliverable. Footnotes (Phase 4), citations (Phase 5),
% and Wade-Giles → Pinyin (Phase 6) are not yet applied here.

\chapter{The Ratio of Scapulas to Plastrons}
\label{app:2}

% source: scan 178, printed 161
% intro-echo-dropped: The Ratio of Scapulas
% intro-echo-dropped: to Plastrons
% intro-echo-dropped: I.
Fragment Counts
\appendixsectionlabel{sec:appendices-app02:1}{1}
New finds may always alter our count of the bones and shells used by the Shang diviners.
Any conclusion about the ratio of plastrons to scapulas, therefore, is approximate and provisional.
The ponderous calculations that follow represent no more than a series of educated guesses.
Many scholars have thought that the Shang used shell more frequently than bone.\footnote[1]{E.g., \pinyinterm{zhang-bingquan} (1954), p. 218.} \pinyinterm{shi-zhangru} argued that the corpus we now have could be recombined to form 8,000-plus plastrons
and 4,000-plus scapulas, that is, in a shell-bone ratio of 2 to 1.\footnote[2]{\pinyinterm{shi-zhangru} (1947), p. 45.} %2
According to \pinyinterm{zhou-hongxiang-tight},
``the ratio of plastrons to extant scapulae is 3 to 1.''\footnote[3]{\pinyinterm{zhou-hongxiang-tight} (1973), p. 177.} %3
\pinyinterm{hu-houxuan}, by contrast, calculated that
the Shang had used 16,003 plastrons and 11,858 scapulas, that is, in an approximate shell-bone
ratio of 4 to 3.\footnote[4]{\pinyinterm{hu-houxuan} (1944a), p. 5b.} %4
Two issues are involved in estimates of this sort. First, what is the ratio between the number
of bone and shell fragments in the present corpus of approximately 176,000 fragments (see appendix
\ref{app:3}, sec. \ref{sec:appendices-app03:1}% was: (appendix 3, sec. 1)
)? This we do not yet know. Short of undertaking a laborious piece-by-piece census, the
best procedure is to calculate the number of bone and shell fragments for the large collections, for
which such counts may readily be made, and assume that due both to their size and the random
nature of their provenance, such collections are representative of the entire corpus. Second, having
arrived at the ratio between the number of bone and shell fragments, we must find out how many
bone fragments would make a scapula and how many shell fragments would make a plastron.\footnote[5]{In what follows, I make no attempt to distinguish plastrons from carapaces, which were used far less frequently (sec.~\ref{sec:chapters-ch01:1.3.2}% was: (sec. 1.3.2)
).}
This again requires a series of assumptions: about the size of the average scapula and plastron and
about the size of the average bone and shell fragment.
To estimate, first, the bone-shell ratio and its variations with time, we may work with large
collections of varied provenance whose editors have recorded whether each fragment is bone or
shell. The ratio of bone and shell fragments in \pinyinterm{jiabian}, \pinyinterm{renwen}, Menzies, and White is given in
table 32. The figures for Menzies and White are not fully comparable with those for \pinyinterm{jiabian} and
\pinyinterm{renwen} since \pinyinterm{xu-jinxiong-tight} places the inscriptions of the Royal Family group in period IV,
% source: scan 179, printed 162
while \pinyinterm{qu-wanli} and Kaizuka and Itō place, or permit us to place, them in period I.\footnote[7]{On this issue, see ch. 2, n. 18.} Nevertheless,
the general totals should have some validity. Similar counts need to be made of other large, randomly
assembled, collections, such as \pinyinterm{jingjin} and \pinyinterm{xucun}.\footnote[8]{For similar tabulations that seek to summarize the situation for the whole corpus of inscriptions, see Tung (1933), p. 420; \pinyinterm{chen-mengjia} (1956), p. 8.}
The \pinyinterm{renwen} and Menzies figures support the view that there was a virtual monopoly of bone
in period III+IV and a marked preference for shell in period V, a phenomenon that has been
well remarked.\footnote[9]{\pinyinterm{chen-mengjia} (1956), p. 8; \pinyinterm{renwen} 1096, shakubun, p. 358, n.; Itō (1962a), p. 253; \pinyinterm{xu-jinxiong-tight} (1973a), pp. 26, 126.} The \pinyinterm{jiabian} figures, however, reveal a slight preference for shell in III+IV, and
more nearly equal usage of bone and shell in V. \pinyinterm{xu-jinxiong-tight} believes that period III inscriptions
were mainly recorded on shell, and that period IV inscriptions, with a few exceptions in IVb.
were all on bone.\footnote[10]{\pinyinterm{xu-jinxiong-tight} (1973a), pp. 2-3, 48-49. But see n. 11.} The extent to which this conclusion is supported by the datings of the editors
of our four collections is indicated by table 5.\footnote[11]{On the problems involved in separating period III and IV inscriptions, see ch. 4, n. 62 and n. 161.} The 1973 finds, largely on bone, and largely, it is
said, from period IV, give added support to this view.\footnote[12]{See ch. 4, n. 188.}
2. Scapula and Plastron Counts
\appendixsectionlabel{sec:appendices-app02:2}{2}
Such counts, however, tell us only about the relative frequency of fragments; they do
not necessarily tell us about the relative frequency of whole scapulas and plastrons. Before drawing
conclusions about frequency of usage, it is necessary to determine the relative size of the bone
and shell fragments, and the relative size of scapulas and plastrons. For ten large pieces of bone
and ten small pieces of shell do not necessarily indicate that scapulas and plastrons were used in
equal proportion.
\pinyinterm{shi-zhangru} argued that it would normally take twenty fragments to make one plastron
or scapula.\footnote[13]{\pinyinterm{shi-zhangru} (1947), p. 45.} \pinyinterm{hu-houxuan} suggested, as a tentative approximation, that although it might take
ten shell fragments to make a plastron, only five bone fragments would be needed to make a scapula.\footnote[14]{\pinyinterm{hu-houxuan} (1944a), p. 5b.}
Hu is undoubtedly correct in assuming that more shell than bone would be required to reconstitute
one piece. Plastrons tended to break into smaller fragments because of their suture lines, which are
not present in scapulas.
To determine the plastron-to-scapula ratio, however, it is first necessary to determine the
size of the average bone and the average shell fragment. By determining the number of fragments
reproduced per plate in large, representative collections, it is possible to select the average plate
(or plates) and then calculate the average surface area of the fragments on that plate (or plates);
for example, \pinyinterm{jiabian} reproduces 2,513 shell fragments on 142.6 plates; there is thus an average
of 18 shell fragments per plate; the surface area of the fragments on plates 8 and 9 (containing a
total of 36 fragments, or 18 fragments per plate), measured by imposing a metric grid over the
rubbings, is found to be 10.6 sq. cms.; this may be taken as the average size of all \pinyinterm{jiabian} shell
fragments. Table 33 indicates the situation in \pinyinterm{jiabian}, \pinyinterm{renwen}, and Menzies, and, for comparison,
in \pinyinterm{yibian} and Tieyun, two nontypical collections.\footnote[6]{I exclude \pinyinterm{yibian} from these discussions. It is composed primarily of period I and RFG shell fragments from one pit, YH 127, and cannot be taken as a representative sample in terms of either chronological range or divination material.} The unusual smallness of the Menzies shell
fragments, already well known and confirmed by the table, suggests that the collection is atypical
% source: scan 180, printed 163
and should be excluded from our calculations.\footnote[15]{On the smallness of the fragments, see \textcite{Mickel1973Review}. Menzies, presumably, could not afford many larger shell fragments; his better examples had already been published in Van-pei, ``Ming,'' and Yinxu.} %15
 Similarly, Tieyun is atypical in that the shell
fragments tend to be larger than the bone fragments-the result presumably of Liu E's selection
of his best and biggest pieces for this first publication of oracle-bone rubbings. \pinyinterm{yibian} should also
be excluded from our calculations of average surface area because its fragments, mostly excavated
from one pit, YH127, and comparatively unbroken, are atypically large (although the ratio
between the size of the bone and shell fragments is typical).
If we are willing to accept \pinyinterm{jiabian} and \pinyinterm{renwen} as representative, we may reach the following
conclusions by averaging the figures given for these two collections in table 33: the average bone
fragment had an area of approximately 19 sq. cms., the average shell fragment had an area of
approximately 9 sq. cms.; for the average fragment, the ratio of bone-to-shell surface area was
approximately 2 to 1.
Before we can translate these figures into scapulas and plastrons, it is still necessary to estimate
the surface area of the average scapula and plastron. This is hard to determine, especially in the
case of the scapulas, for which whole samples are rare. As a result of measuring all the whole or
nearly whole scapulas I could find, I have, faute de mieux, selected 寧滬 1.119 (D) as an average
specimen; it has an area of about 310 sq. cms. The average length (PL) of the first 118 measurable
plastrons (excluding duplicates) in \pinyinterm{pinbian} is 24.6 cms., and the average width of the first 119
measurable plastrons (not always the same as the 118 whose length was measured) is 17.1 cms.
From this it may be estimated that the average area of the plastrons was approximately 318 sq.
cms. Combining these figures with those for the average fragment area given above, we find that
it would have taken approximately 16 bone fragments (310 ÷ 19) to make one scapula, and
approximately 35 shell fragments (318 ÷ 9) to make one plastron. Such figures would not apply
to every collection; in some the fragments are far larger, in others far smaller. Ideally, bone and
shell fragments in every collection should be measured, duplicates excluded, and the average
fragment areas determined for each period. But until such work is done-and I doubt that the
results would be worth the painstaking effort-these initial approximations may be accepted as
being of the right magnitude for calculations concerning the corpus as a whole.
If we translate the figures in table 31 into terms of whole scapulas and plastrons and use the
ratios obtained for each collection (table 33), we arrive at the conclusions given in tables 34 and
35. These are interesting because, crude though they are, they suggest that with the exception of
period III+IV, when bone was more popular, scapulas and plastrons were used in roughly equal
numbers during the historical period.
It may be objected that merely throwing the \pinyinterm{yibian} fragments into the balance would immediately reestablish the dominance of shell over bone. The collection contains over 17,000 shell
fragments found in YH 127, which were accompanied by only 8 bone fragments;\footnote[16]{There is some uncertainty about the exact number. Tung (1949a), p. 259, gives registration numbers for 17,087 fragments; \pinyinterm{shi-zhangru} (1959), p. 322, in apparent agreement with \pinyinterm{chen-mengjia} (1956), p. 156, gives registration numbers for 17,129 fragments.} the shell fragments
would be equivalent to approximately 1,700 plastrons.\footnote[17]{Using the \pinyinterm{yibian} ratio (table 33) of 10.4 fragments per plastron, I estimate that some 800-plus plastrons (or carapaces) have so far been reconstructed, wholly or substantially, in \pinyinterm{yibian} and \pinyinterm{pinbian}.} %17
 The 1973 discovery of 7,040 bone fragments, which would be equivalent to approximately 1,700 scapulas and only 110 shell fragments
(approximately 10.5 plastrons), suggests that including additional collections of this sort would
% source: scan 181, printed 164
not alter the conclusions reached above.\footnote[18]{On this find, see ch. 4, n. 188, above. Again, I use the \pinyinterm{yibian} ratios from table 33.} Until the entire corpus is classified and measured, we
must rest content with extrapolations based, so far as possible, on typical collections, not on large,
atypical caches that were almost exclusively bone or shell.
3. Bone, Shell, and Divination Topic
\appendixsectionlabel{sec:appendices-app02:3}{3}
The conclusion (sec.~\ref{sec:chapters-ch01:1.2.3}% was: (sec. 1.2.3)
) that the diviner's choice of bone or shell appears to have
been unrelated to the matter being divined is provisional, but it accords with the opinion of modern
scholars. \pinyinterm{luo-zhenyu} originally suggested that sacrifice divinations were recorded on turtle shells
and other matters on scapulas; that hunt divinations were recorded on leg bones and that scapulas
were preferred for military topics.\footnote[19]{\pinyinterm{luo-zhenyu} (1914), p. 107b; (1969 ed., ch. 3, p. 64b).} %19
 Initial research suggests that this hypothesis is too strict. Using
the ratios given in table 33, and taking all periods together, I find that the \pinyinterm{renwen} and Menzies
figures reveal a Shang preference for recording sacrifice inscriptions on bone.\footnote[20]{In the equations below, I assume (following the argument in appendix \ref{app:3}, sec. \ref{sec:appendices-app03:1}% was: (appendix 3, sec. 1)
) that inscribed fragments should be combined with uninscribed fragments in the ratio of 52 to 48. \pinyinterm{renwen}: 432 inscribed shell fragments = ca. 21 plastrons; 373 inscribed bone fragments = ca. 37 scapulas. Menzies: 839-plus inscribed shell fragments = ca. 23 plastrons; 309-plus inscribed bone fragments = ca. 48 scapulas.} %20
 The \pinyinterm{renwen} fragments reveal a strong preference for recording hunt inscriptions on bone;\footnote[21]{Fifty-seven inscribed shell fragments = ca. 3 plastrons; 176 inscribed bone fragments = ca. 17 scapulas.} %21
 a similar preference
for bone appears in Menzies.\footnote[22]{Ca. 72 inscribed shell fragments = ca. 2 plastrons; ca. 145 inscribed bone fragments = ca. 17 scapulas.} But the situation in \pinyinterm{jiabian} indicates that the Shang used shell
and bone equally for hunt divinations.\footnote[23]{Eighty inscribed shell fragments = ca. 5 plastrons; 41 inscribed bone fragments = ca. 5 scapulas.} %23
 If we identify the materials used for recording divinations
 about the various cattle sacrifices in \pinyinterm{jiabian}, \pinyinterm{pinbian}, \pinyinterm{yibian}, and \pinyinterm{renwen}, we find that bone
and shell were used in about the same proportion.\footnote[24]{Twenty-six inscribed shell fragments = ca. 1.5 plastrons; 14 inscribed bone fragments = ca. 2 scapulas.} %24

These figures may simply reflect the general predominance of bone or shell in the individual
collections counted, but they do at least suggest that scapulas and plastrons could be used in
significant numbers for divining either sacrifices or hunts; whatever religious or magical considerations may have once related the divination bone to the bone of the animal being sacrificed or
hunted, they were not strongly in evidence by the historical period.\footnote[25]{The theology of Shang divination will be taken up in Studies. On this particular point, cf. Hultkrantz (1968), p. 78.} %25

Firmer conclusions about relative preference for bone or shell would have to be made by the
tedious and difficult classification of every oracle bone, topic by topic, in every collection. The
count would also need to be refined by considering the preferences of particular diviners and
changes in preference in the various periods.\footnote[26]{Cf. ch. 1, n. 34.} %26
